\scanpage{364}
\mspage{88}
que la catégorie $\mathrm{Hot}_{\widetilde{\mathbb W}_\infty}$ mérite d'être appelée catégorie des ``\uline{types d'homotopie simples}'' (à la Whitehead)\textsuperscript{*}.%
% slanted
\margin{\unsure{Canulé|Annulé} (\unsure{?!}), cf commentaire 6.5 page 92.} Un premier test décisif, c'est de vérifier que l'invariant ``\unsure{chronologique|cohomologique} principal de Whitehead d'un type d'\deleted{\unsure{R}}hom. simple, à savoir un élément dans un $K(\mathbb Z(\pi_1))$ (si je me rappelle bien)\textsuperscript{*} -- ou quelque
% slanted
\margin{non, c'est un invariant plus fin\ldots} chose d'approchant, peut se définir déjà en termes \inserted{des objets} de $\mathrm{Hot}_{\bullet\widetilde{\mathbb W}_\infty}$ (le $\mathrm{Hot}$ pointé)\textsuperscript{*}. Et d'abord, si on associe à une \deleted{\unsure{espace} \inserted{\unsure{avec}} \unsure{triang.}} ``maquette simpliciale'' finie (\unsure{décrivant|divisant} \unsure{combinatoirement} une triangulation \unsure{finie}) l'ens. ordonné de ses simplexes, vu comme un objet de $\mathrm{Cat}$, l'objet correspondant de $\mathrm{Hot}_{\widetilde{\mathbb W}_\infty}$, celui-ci ne change pas\textsuperscript{*}, à
% slanted
\margin{C'est évident \unsure{$=$|même} pour $W_{\unsure{\omega|w}}$, puisque le type d'homotopie \uuline{\unsure{strict|stricte}} \uline{ne change} pas ! Non, \unsure{il y a|s'il y a} \unsure{ici} confusion entre homotopie \unsure{topologique} et homotopie \unsure{combinatoire|combinatoires}, p. 92\ldots} isom. can. près, quand on \deleted{\unsure{remplace}} \inserted{effectue} sur \deleted{\unsure{la}} le ``polyèdre'' envisagé une opération élé%
% archive sheet number: 363
